\newcounter{stoppagebreaks}
\documentclass{article}

% packages
\usepackage{amsmath, amsthm, thmtools, amsfonts, amssymb, luacode, catchfile, hyperref, ifthen, graphicx, bbm}
\ifcsname c@kobocompile\endcsname
	\usepackage[a5paper, total={1072pt, 1448pt}, margin=10pt, includeheadfoot]{geometry} % set page margins
\else
	\usepackage[a4paper, margin=50pt, includeheadfoot]{geometry}
\fi
\usepackage[shortlabels]{enumitem}
\usepackage[skip=3pt, indent=0pt]{parskip}

% language
\usepackage[bidi=basic, layout=tabular, provide=*]{babel}
\ifcsname c@english\endcsname
	\babelprovide[main, import]{english}
\else
	\babelprovide[main, import, alph=letters.gershayim]{hebrew}
	\setlist[enumerate,2]{label={\alph*.}}
	\babelprovide{rl}
\fi
\babelfont{rm}[Renderer=Harfbuzz]{Libertinus Serif}
\babelfont{sf}{Libertinus Sans}
\babelfont{tt}{Libertinus Mono}

% style
\ifcsname c@stoppagebreaks\endcsname
\else
	\AddToHook{cmd/section/before}{\clearpage}	% Add line break before section
\fi
\linespread{1.3}
\setcounter{secnumdepth}{0}		% Remove default number tags from sections, this won't do well with theorems
\AtBeginDocument{\setlength{\belowdisplayskip}{3pt}}
\AtBeginDocument{\setlength{\abovedisplayskip}{3pt}}
\graphicspath{ {../images/} }

% operators
\DeclareMathOperator\cis{cis}
\DeclareMathOperator\Sp{Span}
\DeclareMathOperator\tr{tr}
\DeclareMathOperator\im{Im}
\DeclareMathOperator\re{Re}
\DeclareMathOperator\diag{diag}
\DeclareMathOperator*\lowlim{\underline{lim}}
\DeclareMathOperator*\uplim{\overline{lim}}
\DeclareMathOperator\rng{rng}
\DeclareMathOperator\Sym{Sym}
\DeclareMathOperator\Arg{Arg}
\DeclareMathOperator\Log{Log}
\DeclareMathOperator\dom{dom}
\DeclareMathOperator\supp{Supp}
\DeclareMathOperator\var{Var}
\DeclareMathOperator\cov{Cov}
\DeclareMathOperator\End{End}
\DeclareMathOperator\id{id}
\DeclareMathOperator\cl{cl}
\DeclareMathOperator\gal{Gal}
\DeclareMathOperator\dist{dist}
\DeclareMathOperator\form{form}
\DeclareMathOperator\term{term}
\DeclareMathOperator\sent{sent}
\DeclareMathOperator\Th{Th}
\DeclareMathOperator\Hom{Hom}
\DeclareMathOperator\Aut{Aut}
\DeclareMathOperator\Char{char}
\DeclareMathOperator\vol{vol}
\DeclareMathOperator\GL{GL}
\DeclareMathOperator\SL{SL}
\DeclareMathOperator\lcm{lcm}

% commands
\newcommand{\cupdot}{\mathbin{\mathaccent\cdot\cup}}
\newcommand{\indicator}{\mathbbm{1}}
\newcommand{\AAA}[0]{\mathbb{A}}
\newcommand{\Aa}[0]{\mathcal{A}}
\newcommand{\BB}[0]{\mathbb{B}}
\newcommand{\Bb}[0]{\mathcal{B}}
\newcommand{\CC}[0]{\mathbb{C}}
\newcommand{\Cc}[0]{\mathcal{C}}
\newcommand{\DD}[0]{\mathbb{D}}
\newcommand{\Dd}[0]{\mathcal{D}}
\newcommand{\EE}[0]{\mathbb{E}}
\newcommand{\FF}[0]{\mathbb{F}}
\newcommand{\Ff}[0]{\mathcal{F}}
\newcommand{\Ii}[0]{\mathcal{I}}
\newcommand{\Gg}[0]{\mathcal{G}}
\newcommand{\HH}[0]{\mathbb{H}}
\newcommand{\Hh}[0]{\mathcal{H}}
\newcommand{\Kk}[0]{\mathcal{K}}
\newcommand{\Ll}[0]{\mathcal{L}}
\newcommand{\Mm}[0]{\mathcal{M}}
\newcommand{\NN}[0]{\mathbb{N}}
\newcommand{\Nn}[0]{\mathcal{N}}
\newcommand{\PP}[0]{\mathbb{P}}
\newcommand{\Pp}[0]{\mathcal{P}}
\newcommand{\QQ}[0]{\mathbb{Q}}
\newcommand{\RR}[0]{\mathbb{R}}
\newcommand{\Rr}[0]{\mathcal{R}}
\newcommand{\Ss}[0]{\mathcal{S}}
\newcommand{\TT}[0]{\mathbb{T}}
\newcommand{\Tt}[0]{\mathcal{T}}
\newcommand{\UU}[0]{\mathbb{U}}
\newcommand{\Uu}[0]{\mathcal{U}}
\newcommand{\Vv}[0]{\mathcal{V}}
\newcommand{\Ww}[0]{\mathcal{W}}
\newcommand{\Xx}[0]{\mathcal{X}}
\newcommand{\Yy}[0]{\mathcal{Y}}
\newcommand{\ZZ}[0]{\mathbb{Z}}
\newcommand{\Zz}[0]{\mathcal{Z}}
\newcommand{\acts}[0]{\circlearrowright}
\newcommand{\explain}[2] {
	\begin{flalign*}
		 && \text{#2} && \text{#1}
	\end{flalign*}
}
\newcommand{\maketitleprint}[1][green]{
	\begin{center}
		\includegraphics[width=6cm]{dragon_#1}
	\end{center}
}

% theorem commands
\newtheoremstyle{c_remark}
	{}	% Space above
	{}	% Space below
	{}% Body font
	{}	% Indent amount
	{\bfseries}	% Theorem head font
	{}	% Punctuation after theorem head
	{.5em}	% Space after theorem head
	{\thmname{#1}\thmnumber{ #2}\thmnote{ \normalfont{\text{(#3)}}}}	% head content
\newtheoremstyle{c_definition}
	{3pt}	% Space above
	{3pt}	% Space below
	{}% Body font
	{}	% Indent amount
	{\bfseries}	% Theorem head font
	{}	% Punctuation after theorem head
	{.5em}	% Space after theorem head
	{\thmname{#1}\thmnumber{ #2}\thmnote{ \normalfont{\text{(#3)}}}}	% head content
\newtheoremstyle{c_plain}
	{3pt}	% Space above
	{3pt}	% Space below
	{\itshape}% Body font
	{}	% Indent amount
	{\bfseries}	% Theorem head font
	{}	% Punctuation after theorem head
	{.5em}	% Space after theorem head
	{\thmname{#1}\thmnumber{ #2}\thmnote{ \text{(#3)}}}	% head content

\ifcsname c@english\endcsname
	\theoremstyle{plain}
	\newtheorem{theorem}{Theorem}[section]
	\newtheorem*{theorem*}{Theorem}
	\newtheorem{lemma}[theorem]{Lemma}
	\newtheorem{proposition}[theorem]{Proposition}
	\newtheorem*{proposition*}{Proposition}

	\theoremstyle{definition}
	\newtheorem{definition}[theorem]{Definition}
	\newtheorem*{definition*}{Definition}
	\newtheorem{example}{Example}[section]
	\newtheorem{exercise}{Exercise}[section]

	\theoremstyle{remark}
	\newtheorem*{remark}{Remark}
	\newtheorem*{solution}{Solution}
	\newtheorem{corollary}[theorem]{Corollary}
	\newtheorem{notation}[theorem]{Notation}
\else
	\theoremstyle{c_plain}
	\newtheorem{theorem}{משפט}[section]
	\newtheorem*{theorem*}{משפט}
	\newtheorem{lemma}[theorem]{למה}
	\newtheorem{proposition}[theorem]{טענה}
	\newtheorem*{proposition*}{טענה}

	\theoremstyle{c_definition}
	\newtheorem{definition}[theorem]{הגדרה}
	\newtheorem*{definition*}{הגדרה}
	\newtheorem{example}{דוגמה}[section]
	\newtheorem{exercise}{תרגיל}[section]

	\theoremstyle{c_remark}
	\newtheorem*{remark}{הערה}
	\newtheorem*{solution}{פתרון}
	\newtheorem{corollary}[theorem]{מסקנה}
	\newtheorem{notation}[theorem]{סימון}
\fi

% Questions related commands
\newcounter{question}
\setcounter{question}{1}
\newcounter{sub_question}
\setcounter{sub_question}{1}

\ifcsname c@english\endcsname
	\newcommand{\question}[1][0]{
		\ifthenelse{#1 = 0}{}{\setcounter{question}{#1}}
		\section{Question \arabic{question}}
		\addtocounter{question}{1}
		\setcounter{sub_question}{1}
	}

	\newcommand{\subquestion}[1][0]{
		\ifthenelse{#1 = 0}{}{\setcounter{sub_question}{#1}}
		\subsection{Part \alph{sub_question}}
		\addtocounter{sub_question}{1}
	}
\else
	\newcommand{\question}[1][0]{
		\ifthenelse{#1 = 0}{}{\setcounter{question}{#1}}
		\section{שאלה \arabic{question}}
		\addtocounter{question}{1}
		\setcounter{sub_question}{1}
	}

	\newcommand{\subquestion}[1][0]{
		\ifthenelse{#1 = 0}{}{\setcounter{sub_question}{#1}}
		\subsection{סעיף \localecounter{letters.gershayim}{sub_question}}
		\addtocounter{sub_question}{1}
	}
\fi

% import lua and start of document
\directlua{common = require ('../common')}

\GetEnv{AUTHOR}

% headers
\author{\AUTHOR}
\date\today

\setcounter{secnumdepth}{2}
% chktex-file 9
% chktex-file 17

\usepackage{pgfplots}
\pgfplotsset{compat=1.18}

\title{מתמטיקה א' --- תרגול 6}
\date{}

\begin{document}

\maketitle

\section{חסימות}
\begin{definition}[פונקציה חסומה]
	נאמר שפונקציה $f : A \to B$ היא חסומה אם קיימים $a, b \in \RR$ כך שלכל $x$ מתקיים,
	\[
		a \le f(x) \le b
	\]
\end{definition}
באופן מאוד דומה נוכל להגיד שהיא חסומה מלמטה אם קיים רק $a$ כזה, או חסומה מלמעלה אם יש רק $b$ כזה.
\begin{example}
	נראה כי $f(x) = \frac{1}{x^2 + 1}$ חסומה:

	$x^2$ תמיד אי־שלילי, ולכן גם $x^2 + 1$ חיובית ממש, וחלוקה של חיובי ($1$) וחיובי ($x^2 + 1$) היא גם חיובית, לכן $0 \le f(x)$ תמיד, כלומר $f$ חסומה מלמטה. \\
	כדי להראות שהיא חסומה מלמעלה נשים לב שמתקיים,
	\[
		1 \le x^2 + 1
		\iff \frac{1}{1} \ge \frac{1}{x^2 + 1}
	\]
	כלומר $f(x) \le 1$ לכל $x$, ונסיק שהפונקציה חסומה.
\end{example}
נראה דוגמה לפונקציה לא חסומה.
\begin{example}
	נגדיר את $f(x) = \exp(x)$, אז אנחנו כבר יודעים ש־$f : \RR \to (0, \infty)$, כלומר הפונקציה תמיד מחזירה מספר חיובי, ולכן היא חסומה מלמטה על־ידי $0$. \\
	מהצד השני, נרצה להראות שהיא לא חסומה מלמעלה, כדי לעשות זאת נבצע הוכחה על דרך השלילה, כלומר נראה שאילו הייתה הפונקציה חסומה מלמעלה, אז היינו מקבלים סתירה מתמטית.
	נניח בשלילה שהיא חסומה מלמעלה, כלומר שלכל $x$ מתקיים $f(x) \le b$ לאיזשהו מספר $b \in \RR$.
	לכן $\exp(x) \le b$, אבל מהפיכות של $\exp$ ושל $\ln$ נקבל את אי־השוויון $x \le \ln b$, ואנחנו יודעים שלכל $b$ שנבחר תמיד יש $x$ קטן מספיק ככה שאי־השוויון מתקיים, לדוגמה $x = \ln(b) - 1$.
	נסיק אם כך שהטענה $\exp(x) \le b$ היא שקר, כלומר סתירה, ולכן נסיק שזה לא יתכן, הפונקציה $f$ לא חסומה מלמעלה ובפרט לא חסומה בכלל.
\end{example}
\begin{exercise}
	ננסה לנחש ללא חישוב אם כל אחת מהפונקציות הבאות חסומה (מלמטה, מלמעלה או שניהם):
	\begin{enumerate}
		\item $f(x) = x$
		\item $f : [0, 1] \to \RR$ המוגדרת על־ידי $f(x) = x$
		\item $f(x) = x^2$
		\item $f(x) = \frac{1}{x}$
	\end{enumerate}
\end{exercise}
\begin{solution}
	\begin{enumerate}
		\item הפונקציה לא חסומה
		\item הפונקציה כן חסומה, כי $0 \le f(x) \le 1$ לכל מספר שנציב מהתחום
		\item הפונקציה חסומה מלמטה, הרי ריבוע של כל מספר הוא חיובי, אבל היא לא חסומה מלמעלה
		\item הפונקציה לא חסומה משום צד
	\end{enumerate}
\end{solution}

\section{זוגיות ואי־זוגיות}
\begin{definition}[פונקציה זוגית ואי־זוגית]
	פונקציה $f : A \to B$ נקראת זוגית אם,
	\[
		\forall x \in A: f(-x) = f(x)
	\]
	והיא נקראת אי־זוגית אם,
	\[
		\forall x \in A: f(-x) = -f(x)
	\]
\end{definition}
נוכל להגיד שהפונקציה היא זוגית אם שני החצאים שלה סביב ציר ה־$y$ הם סימטריים, באופן שקול שהגרף שלה מהווה מראה בין החצי הימני והשמאלי.
מהצד השני נוכל להגיד שפונקציה היא אי־זוגית אם אפשר לסובב את הגרף של הפונקציה ב־$180$ מעלות ולקבל את אותו הגרף.
\begin{example}
	נגדיר את הפונקציה $f(x) = x^2$, אז $f$ היא פונקציה זוגית.
	לכל $x$ מתקיים $f(x) = x^2$ וגם $f(-x) = {(-x)}^2 = x^2$, ולכן $f(x) = f(-x)$.
\end{example}
\begin{example}
	הפונקציה $f(x) = x^3$ היא אי־זוגית.
	לכל $x$ מתקיים $f(-x) = {(-x)}^3 = {(-1)}^3 \cdot x^3 = -x^3 = -f(x)$.
\end{example}
\begin{example}
	הפונקציה $f(x) = x + 1$ היא לא זוגית ולא אי־זוגית.
	נבדוק זאת על־ידי הצבה (זכרו שכדי לבדוק שמשהו לא נכון מספיק למצוא ערך אחד עבורו התנאי לא מתקיים):
	\[
		f(3)
		= 3 + 1
		= 4
		\ne -2
		= -3 + 1
		= f(-3)
	\]
	כלומר על־ידי הצבת $x = 3$ נראה שהיא לא זוגית, ועבור אי־זוגיות,
	\[
		f(-3)
		= -2
		\qquad
		-f(3)
		= -4
	\]
	ולכן מצאנו שהיא גם לא אי־זוגית.
\end{example}
\begin{exercise}
	נתונה לנו הפונקציה $f : [0, \infty) \to \RR$ המוגדרת על־ידי $f(x) = x$. \\
	השלימו את $f$ להיות פונקציה זוגית על־ידי הגדרת הערכים שלה למספרים השליליים.
\end{exercise}
\begin{solution}
	המטרה שלנו היא למצוא פונקציה חדשה $g : \RR \to \RR$ ככה שלכל מספר אי־שלילי $x$ מתקיים,
	\[
		g(x)
		= f(x)
	\]
	ולכל מספר שלילי מתקיים,
	\[
		g(x)
		= g(-x)
		= f(-x)
		= -x
	\]
	כלומר נוכל להגדיר את,
	\[
		g(x)
		= \begin{cases}
			x & x \ge 0 \\
			-x & x < 0
		\end{cases}
	\]
	זוהי למעשה בדיוק פונקציית הערך המוחלט! כלומר $|x|$ היא פונקציה זוגית.
\end{solution}

\section{פונקציות טריגונומטריות}
בהרצאה ראינו איך מגדירים את הפונקציות הטריגונומטריות $\sin, \cos, \tan$ על־ידי שימוש במעגל היחידה, ודיברנו על חישובים שלהן.
\begin{exercise}
	המירו לרדיאנים את הזוויות: $0^\circ, 60^\circ, 90^\circ$.
\end{exercise}
\begin{solution}
	$0$ מעלות הן בדיוק $0$ רדיאנים, זאת שכן הזווית ''מנוונת''. \\
	$60^\circ$ הן בדיוק שישית הסיבוב, וסיבוב שלם ברדיאנים הוא $2 \pi$, לכן נסיק שמתקיים $60^\circ = \frac{2 \pi}{6} = \frac{\pi}{3}$. \\
	באופן דומה $90^\circ$ הן בדיוק רבע הסיבוב, ולכן בהמרה לרדיאנים נקבל $\frac{2 \pi}{4} = \frac{\pi}{2}$.
\end{solution}
ניזכר כי בהרצאה ראינו שסינוס וקוסינוס הן פונקציות אי־זוגית וזוגית בהתאמה, שתיהן חסומות ומקבלות ערכים $[-1, 1]$ בלבד. \\
נעבור לפתרון משוואה טריגונומטרית.
\begin{exercise}
	פתרו את המשוואה:
	\[
		\cos(x)
		= 1
	\]
\end{exercise}
\begin{solution}
	נמיר את השאלה לשאלה על מעגל היחידה, מתי ערך ה־$x$ הנוצר על־ידי זווית מגיע בדיוק ל־$1$:
	\begin{otherlanguage}{english}
		\begin{center}
			\begin{tikzpicture}
				\begin{axis}[
					axis lines=middle,
					axis equal image,
					xmin=-1.4, xmax=1.4,
					ymin=-1.2, ymax=1.2,
					xlabel={$x$},
					ylabel={$y$},
					ticks=none,
					clip=false
					]
					\addplot[blue, thick, domain=0:360, samples=200]
					({cos(x)}, {sin(x)});

					\addplot[red, thick, domain=-1.1:1.1] ({1}, {x});
					\node[red, right] at (axis cs:1,0.75) {$x=1$};

					\addplot[only marks, mark=*, mark size=2pt] coordinates {(1,0)};
					\node[below right] at (axis cs:1,0) {$(1,0)$};
				\end{axis}
			\end{tikzpicture}
		\end{center}
	\end{otherlanguage}
	נשים לב שהזווית מגיעה בציר ה־$x$ ל־$1$ רק כשעושים מספר שלם של סיבובים, כלומר אם $x \in \{ 2 \pi k \mid k \in \ZZ \}$.
\end{solution}

\section{מחזוריות}
\begin{definition}[פונקציה מחזורית]
	פונקציה $f : A \to B$ נקראת $L$־מחזורית אם לכל $x, x + L \in A$ מתקיים $f(x) = f(x + L)$.
\end{definition}
ראינו בהרצאה כי $\sin, \cos, \tan$ הן מחזוריות.
\begin{exercise}
	הראו כי ${(\sin(2 x))}^2$ היא מחזורית.
\end{exercise}
\begin{solution}
	אנו יודעים כי $\sin x$ היא $2 \pi$־מחזורית, כלומר,
	\[
		\sin(x)
		= \sin(x + 2 \pi)
	\]
	לכל $x \in \RR$, אם נציב $x = 2 x_1$ נקבל,
	\[
		\sin(2 x_1)
		= \sin(2 x_1 + 2 \pi)
		= \sin(2 (x_1 + \pi))
	\]
	ולכן נסיק ש־$\sin(2 x)$ היא $\pi$־מחזורית.
	נשאר לבדוק את הריבוע,
	\[
		{(\sin(2 x))}^2
		= \sin(2 x) \cdot \sin(2 x)
		= \sin(2 (x + \pi)) \cdot \sin(2 (x + \pi))
		= {(\sin(2 (x + \pi)))}^2
	\]
	כלומר הפונקציה היא אכן $\pi$־מחזורית.
\end{solution}
\begin{exercise}
	תהי פונקציה $f : \RR \to \RR$ המקיימת את התכונות הבאות:
	\begin{enumerate}
		\item עבור $x \in [-2, 2]$ מתקיים $f(x) = x^2$
		\item $f$ היא $4$־מחזורית
	\end{enumerate}
	חשבו את $f(0), f(5), f(31), f(1522)$.
\end{exercise}
\begin{solution}
	$0 \in [-2, 2]$ ולכן נסיק ש־$f(0) = 0^2 = 0$. \\
	ממחזוריות נקבל $f(5) = f(5 - 4) = f(1)$, אבל שוב $1 \in [-2, 2]$ ולכן $f(5) = 1^2 = 1$. \\
	נשתמש בשיטה דומה עבור $31, 1522$:
	\[
		f(31)
		= f(31 - 32)
		= f(-1)
		= {(-1)}^2
		= 1,
		\qquad
		f(1522)
		= f(1522 - 1524)
		= f(-2)
		= 4
	\]
\end{solution}
\begin{exercise}
	בשאלה קודמת קיבלנו פונקציה המוגדרת על חלק מהציר הממשי והרחבנו אותה לפונקציה סימטרית.
	הפעם נניח ש־$f : [1, 2) \to \RR$ מוגדרת על־ידי $f(x) = \frac{1}{x}$, תארו פונקציה $g : \RR \to \RR$ שהיא $1$־מחזורית, שווה ל־$f$ בכל הערכים בהם שתי הפונקציות מוגדרות, וחשבו את $g(\frac{33}{2})$.
\end{exercise}
\begin{solution}
	נגדיר את $g$ כך ש־$g(x) = \frac{1}{x}$ עבור $x \in [1, 2)$, וכן $g(x + 1) = g(x), g(x + 2) = g(x), \ldots$.
	נוכל לסמן זאת על־ידי,
	\[
		\forall x \in [0, 1), k \in \ZZ:
		g(x + k)
		= \frac{1}{x}
	\]
	נקבל בצורה דומה לשאלה הקודמת,
	\[
		g(\frac{33}{2})
		= g(\frac{33}{2} - 15)
		= g(\frac{3}{2})
		= \frac{2}{3}
	\]
\end{solution}

\section{מחסנית}
\begin{exercise}
	בדקו האם הפונקציה $f(x) = \sin(3x) \cdot \cos(5x) \cdot \frac{1}{\sin(x^2 + 3)}$ זוגית או אי־זוגית.
\end{exercise}
\begin{solution}
	נבדוק על־ידי הצבת $-x$:
	\[
		f(-x)
		= \sin(-3x) \cdot \cos(-5x) \cdot \frac{1}{\sin({(-x)}^2 + 3)}
		= \sin(3x) \cdot (-\cos(5x)) \cdot \frac{1}{\sin(x^2 + 3)}
		= -f(x)
	\]
	נסיק ש־$f$ היא פונקציה אי־זוגית.
\end{solution}
\begin{exercise}
	חשבו ללא מחשבון את הביטויים הבאים:
	\begin{enumerate}
		\item $\sin \frac{\pi}{4}$
		\item $\sin(80 \frac{1}{4} \pi)$
		\item $\sin(\frac{101 \pi}{2})$
	\end{enumerate}
\end{exercise}
\begin{solution}
	\begin{enumerate}
		\item זוהי זווית של $45^\circ$ מעלות, על־ידי שימוש במשפט פיתגורס על מעגל היחידה, והעובדה שנוצר משולש שווה שוקיים נקבל,
			\begin{otherlanguage}{english}
				\begin{center}
					\begin{tikzpicture}
						\begin{axis}[
							axis lines=middle,
							axis equal image,
							xmin=-1.2, xmax=1.2,
							ymin=-1.2, ymax=1.2,
							xlabel={$x$},
							ylabel={$y$},
							ticks=none,
							clip=false
							]
							\pgfmathsetmacro{\a}{sqrt(2)/2}

							\addplot[gray, thick, domain=0:360, samples=200]
							({cos(x)}, {sin(x)});

							\addplot[blue, thick] coordinates {
								(0,0) (\a,0) (\a,\a) (0,0)
							};

							\addplot[thick] coordinates {
								({\a-0.08},0) ({\a-0.08},0.08) (\a,0.08)
							};

							\node[below] at (axis cs:{\a/2},0) {$\frac{\sqrt{2}}{2}$};
							\node[right] at (axis cs:\a,{\a/2}) {$\frac{\sqrt{2}}{2}$};
							\node[above left] at (axis cs:{\a/2},{\a/2}) {$1$};
							\node at (axis cs:0.2,0.08) {$\frac{\pi}{4}$};

							\addplot[only marks, mark=*, mark size=1.8pt]
							coordinates {(\a,\a)};
						\end{axis}
					\end{tikzpicture}
				\end{center}
			\end{otherlanguage}
			נקבל בדיוק את המשוואה:
			\[
				x^2 + x^2 = 1
				\implies x = \sqrt{\frac{1}{2}} = \frac{1}{\sqrt{2}}
			\]
		\item נשתמש ב־$2 \pi$־מחזוריות של סינוס:
			\[
				\sin(80 \frac{1}{4} \pi)
				= \sin(80 \frac{1}{4} \pi - 40 \cdot 2 \pi)
				= \sin \frac{\pi}{4}
				= \frac{1}{\sqrt{2}}
			\]
		\item שוב נשתמש בתכונת המחזוריות,
			\[
				\sin(\frac{101 \pi}{2})
				= \sin(\frac{101 \pi}{2} - 25 \cdot 2 \pi)
				= \sin(\frac{\pi}{2})
				= 1
			\]
	\end{enumerate}
\end{solution}
\begin{exercise}
	תהי פונקציה $f : \RR \to \RR$ שהיא גם מונוטונית עולה וגם $3$־מחזורית.
	מצאו ביטוי מפורש לפונקציה.
\end{exercise}
\begin{solution}
	נתון לנו כי $f$ היא מחזורית, כלומר $f(x) = f(x + 3)$.
	בנוסף נתון לנו שהיא עולה, ולכן לכל $x < x_1 < x + 3$ נקבל גם $f(x) \le f(x_1) \le f(x + 3) = f(x)$, כלומר $f(x_1) \le f(x)$ וגם $f(x_1) \ge f(x)$, לכן $f(x) = f(x_1)$ לכל $x_1$ כזה.
	נסמן $f(0) = c$ עבור פרמטר כלשהו $c \in \RR$, אז נקבל $f(x) = c$ לכל $x \in \RR$ בלבד, כלומר הפונקציה חייבת להיות קבועה.
\end{solution}

\end{document}
